\documentclass[a4paper,12pt,answers]{exam}
\usepackage[T1]{fontenc}
\usepackage{lmodern}
\usepackage{comment}
\usepackage{color}
\usepackage{fancyvrb}
\usepackage{graphicx}
\usepackage{geometry}
\usepackage{draftwatermark}

\usepackage{amssymb,amsmath}
\usepackage{ifxetex,ifluatex}
\usepackage{fixltx2e} % provides \textsubscript
% use upquote if available, for straight quotes in verbatim environments
\IfFileExists{upquote.sty}{\usepackage{upquote}}{}
\ifnum 0\ifxetex 1\fi\ifluatex 1\fi=0 % if pdftex
  \usepackage[utf8]{inputenc}
\else % if luatex or xelatex
  \ifxetex
    \usepackage{mathspec}
    \usepackage{xltxtra,xunicode}
  \else
    \usepackage{fontspec}
  \fi
  \defaultfontfeatures{Mapping=tex-text,Scale=MatchLowercase}
  \newcommand{\euro}{€}
\fi
% use microtype if available
\IfFileExists{microtype.sty}{\usepackage{microtype}}{}
\usepackage{longtable,booktabs}
\usepackage{graphicx}
% Redefine \includegraphics so that, unless explicit options are
% given, the image width will not exceed the width of the page.
% Images get their normal width if they fit onto the page, but
% are scaled down if they would overflow the margins.
\makeatletter
\def\ScaleIfNeeded{%
  \ifdim\Gin@nat@width>\linewidth
    \linewidth
  \else
    \Gin@nat@width
  \fi
}
\makeatother
\let\Oldincludegraphics\includegraphics
{%
 \catcode`\@=11\relax%
 \gdef\includegraphics{\@ifnextchar[{\Oldincludegraphics}{\Oldincludegraphics[width=\ScaleIfNeeded]}}%
}%
\ifxetex
  \usepackage[setpagesize=false, % page size defined by xetex
              unicode=false, % unicode breaks when used with xetex
              xetex]{hyperref}
\else
  \usepackage[unicode=true]{hyperref}
\fi
\hypersetup{breaklinks=true,
            bookmarks=true,
            pdfauthor={},
            pdftitle={Worksheet: Streams},
            colorlinks=true,
            citecolor=blue,
            urlcolor=blue,
            linkcolor=magenta,
            pdfborder={0 0 0}}
\urlstyle{same}  % don't use monospace font for urls
\setlength{\parindent}{0pt}
\setlength{\parskip}{6pt plus 2pt minus 1pt}
\setlength{\emergencystretch}{3em}  % prevent overfull lines
%\setcounter{secnumdepth}{0}

\geometry{
  %includeheadfoot,
  margin=2.54cm
}

\newcommand{\codeDir}{src}
\newcommand{\tit}{Exercises --- Week Eight}

\title{Exercises --- Week Eight}
\ifprintanswers
\title{\tit \ with outline solutions}
\else
\title{\tit}
\fi
\date{Spring term 2016}
\author{The functional circus continues\ldots}

\renewcommand{\solutiontitle}{\noindent\textbf{Outline Solution:}\par\noindent}
\shadedsolutions

\SetWatermarkLightness{0.9}
\SetWatermarkScale{0.8}
\ifprintanswers
\SetWatermarkText{Outline solutions}
\else
\SetWatermarkText{}
\fi

\begin{document}
\maketitle

You need to have installed the current Scala distribution and the current Java distribution
before commencing these exercises.
Don't forget the documentation for both distributions as they will come in very useful.

For some of these exercises you will need the test code you used in a previous exercise sheet,
namely the classes below the \verb!atomicscala! folder on the repo. Any source code examples
are available under the \verb!week08! folder.

There are four broad areas begin examined:
\begin{enumerate}
\item Generics,
\item Consolidation of existing concepts,
\item Sequences, and	
\item Higher-Order Functions.
\end{enumerate}

\begin{questions}

\section*{Generics}

\question % 5.1.3.2 - ES
Given the following type:
\begin{Verbatim}
sealed trait IntList
final case object End extends IntList
final case class Pair(head: Int, tail: IntList) extends IntList	
\end{Verbatim}
\begin{parts}
\part
Change the name to \verb+LinkedList+ and make the type of data stored 
in the list generic.
\begin{solution}
\VerbatimInput{src/ll/LinkedList.scala}
\end{solution}

\part
Now implement \verb+length+, returning the length of the \verb+LinkedList+. 
Some test cases are below:
\begin{Verbatim}
val example = Pair(1, Pair(2, Pair(3, End()))) 
assert(example.length == 3) 
assert(example.tail.length == 2) 
assert(End().length == 0)	
\end{Verbatim}
\begin{solution}
\VerbatimInput{src/lltwo/LinkedList.scala}	
\end{solution}

\part
Now implement a method \verb!contains! that determines whether or not a 
given item is in the list. Ensure your code works with the following test cases:
\begin{Verbatim}
val example = Pair(1, Pair(2, Pair(3, Empty()))) 
assert(example.contains(3) == true) 
assert(example.contains(4) == false) 
assert(Empty().contains(0) == false)
// This should not compile
// example.contains("not an Int")	
\end{Verbatim}
\begin{solution}
\VerbatimInput{src/llthree/LinkedList.scala}	
\end{solution}

\part
Implement a method \verb!apply! that returns the $n$th item in the list.
Ensure your solu on works with the following test cases:
\begin{Verbatim}
val example = Pair(1, Pair(2, Pair(3, Empty()))) 
assert(example(0) == 1)
assert(example(1) == 2)
assert(example(2) == 3)
assert(try { 
	example(3) 
	false
} catch {
	case e: Exception => true
})	
\end{Verbatim}
[Yes, Scala does have java type exceptions --- although we try to avoid them.]

\begin{solution}
\VerbatimInput{src/llfour/LinkedList.scala}	
\end{solution}

\part
Throwing an exception isn’t very clean. 
Whenever we throw an exception we lose type safety as there is nothing in the 
type system that will remind us to deal with the error. 
It would be much better to return some kind of result that encodes we can 
\emph{succeed} or \emph{fail}. Consider the following type:
\begin{Verbatim}
sealed trait Result[A]
case class Success[A](result: A) extends Result[A]
case class Failure[A](reason: String) extends Result[A]	
\end{Verbatim}
Change your \verb!apply! method so it returns a \verb!Result!, with a failure case indicating 
what went wrong. Here are some test cases to help you:
\begin{Verbatim}
assert(example(0) == Success(1))
assert(example(1) == Success(2))
assert(example(2) == Success(3))
assert(example(3) == Failure("Index out of bounds"))	
\end{Verbatim}
\begin{solution}
\VerbatimInput{src/llfive/LinkedList.scala}	
\end{solution}

\end{parts}

\question % 5.2.3.1 ES
Consider the following code:
\VerbatimInput{src/IntList.scala}
All of these methods have the same general pattern, which is not surprising as they all 
use structural recursion. It would be nice to be able to remove the duplication.
Below we have sketched out an abstraction over \verb!sum!, 
\verb!length!, and \verb!product!
\begin{Verbatim}
def abstraction(end: Int, f: ???): Int = 
	this match {
		case End => end
		case Pair(hd, tl) => f(hd, tl.abstraction(f, end)) 
	}	
\end{Verbatim}
\begin{parts}
\part
Rename this function to \verb!fold!, which is the name it is usually known as, 
and complete the implementation.

\begin{solution}
\begin{Verbatim}
def fold(end: Int, f: (Int, Int) => Int): Int = this match {
  case End => end
  case Pair(hd, tl) => f(hd, tl.fold(f, end)) 
}	
\end{Verbatim}
\end{solution}

\part 	
Now reimplement \verb!sum!, \verb!length!, and \verb!product! in terms of \verb!fold!.

\begin{solution}
\begin{Verbatim}
sealed trait IntList {
  def fold(end: Int, f: (Int, Int) => Int): Int =
    this match {
      case End => end
      case Pair(hd, tl) => f(hd, tl.fold(end, f))
    }

  def length: Int =
    fold(0, (_, tl) => 1 + tl) 
  
  def product: Int =
    fold(1, (hd, tl) => hd * tl) 
  
  def sum: Int =
    fold(0, (hd, tl) => hd + tl)
}
final case object End extends IntList
final case class Pair(head: Int, tail: IntList) 
  extends IntList	
\end{Verbatim}
\end{solution}

\part 
Is it more convenient to rewrite methods in terms of \verb!fold! 
if they were implemented using pattern matching or polymorphism? 
What does this tell us about the best use of \verb!fold!?

\begin{solution}
When using \verb|fold| in polymorphic implementations we have a lot of duplication; 
the polymorphic implementations without \verb!fold! were simpler to write. 
The pattern matching implementations benefited from \verb!fold! as we remove the 
duplication in the pattern matching.

In general \verb!fold! makes a good interface for users outside the class, 
but not necessarily for use inside the class.
\end{solution}

\part
Why can’t we write our \verb!double! method in terms of \verb!fold!? 
Is it feasible we could if we made some change to \verb!fold!?

\begin{solution}
The types tell us it won’t work. \verb!fold! returns an \verb!Int! and \verb!double! 
returns an \verb!IntList!. However the general structure of \verb!double! is 
captured by \verb!fold!. 
\end{solution}

\part
Implement a generalised version of \verb!fold! and rewrite \verb!double! in terms of it.

\begin{solution}
We want to generalise the return type of \verb!fold!. Our starting point is:
\begin{Verbatim}
def fold(end: Int, f: (Int, Int) => Int): Int	
\end{Verbatim}
Replacing the return type and tracing it back we arrive at:
\begin{Verbatim}
def fold[A](list: IntList, f: (Int, A) => A, end: A): A	
\end{Verbatim}
where we’ve used a generic type on the method to capture the changing return type. 
With this we can implement \verb!double!. When we try to do so we’ll see that type 
inference fails, so we have to give it a bit of help.
\begin{Verbatim}
sealed trait IntList {
  def fold[A](end: A, f: (Int, A) => A): A =
    this match {
      case End => end
      case Pair(hd, tl) => f(hd, tl.fold(end, f))
    }

  def length: Int =
    fold[Int](0, (_, tl) => 1 + tl) 

  def product: Int =
    fold[Int](1, (hd, tl) => hd * tl) 
  
  def sum: Int =
    fold[Int](0, (hd, tl) => hd + tl) 
  
  def double: IntList =
    fold[IntList](End, (hd, tl) => Pair(hd * 2, tl))
}
final case object End extends IntList
final case class Pair(head: Int, tail: IntList) extends IntList	
\end{Verbatim}
\end{solution}
\end{parts}

\section*{Consolidation}

\question
Given the following code:

\VerbatimInput{\codeDir/One.scala}

The \verb!isGoodEnough! test is not very precise for small numbers and might 
lead to non-termination (why?). 
Design a different version of \verb!isGoodEnough! which does not have these problems.	
\begin{solution}
While the number of real numbers is infinite, the number of binary representable numbers 
is finite --– and so is the the number of floating-point numbers that the computer uses.

Floating point numbers can represent most real numbers by approximation only because there 
is no exact binary representation. The ``gap'', or distance, between adjacent floating-point
numbers is smaller for small numbers, and becomes larger as the numbers increase 
(the larger its so-called exponent becomes, actually). Thus, an absolute value 
(a distance of 0.001 in the above code snippet) cannot be used in general-purpose functions.

The solution is to use the \verb!scala.math.ulp(Double): Double! method –--
\verb!ulp! meaning, \emph{units in the last place}, representing 
\emph{the positive distance between this floating-point value and the double 
value next larger in magnitude}.

\VerbatimInput{\codeDir/OneSol.scala}	
\end{solution}

\question
Given the following function:
\begin{Verbatim}
def factorial(n: Int): Int = if (n == 0) 1 else n * factorial(n - 1)
\end{Verbatim}
Design a tail-recursive version of \verb!factorial!.

\begin{solution}
In contrast to other recursive functions, tail-recursive functions can have better performance
as subsequent, recursive, function calls can re-use the memory stack frame that has been reserved
for the function. A function is tail-recursive when the 
call to itself is its very last operation.

\VerbatimInput{\codeDir/TailRec.scala}	
\end{solution}

\question
Can you write a tail-recursive version of a function \verb+sum+, that computes the sum 
of the values of a function at points over a given range, by filling in the \verb+???+’s?
\begin{samepage}
\begin{Verbatim}
def sum(f: Int => Int)(a: Int, b: Int): Int = {
        def iter(a: Int, result: Int): Int = {
                if (???) ???
                else iter(???, ???)
        }
        iter(??, ??)
}	
\end{Verbatim}
\end{samepage}

\begin{solution}
\VerbatimInput{\codeDir/TailSum.scala}	
\end{solution}

\question
Write a function \verb!product! that computes the product of the values of a
function at points over a given range.

\begin{solution}
\VerbatimInput{\codeDir/TailProduct.scala}	
\end{solution}

\question
Write \verb!factorial! in terms of \verb!product!.

\begin{solution}
\ldots	
\end{solution}

\question
Can you write a more general function which generalises both \verb!sum! and \verb!product!?

\begin{solution}
\VerbatimInput{\codeDir/Operate.scala}	
\end{solution}

\question
Given the following code:
\VerbatimInput{\codeDir/IntSet.scala}

\begin{samepage}
Write methods \verb!union! and \verb!intersection! to form the 
union and intersection between two sets.

Add a method
\begin{Verbatim}
def excl(x: Int)
\end{Verbatim}
to return the given set without the element \verb!x!. 
To accomplish this, it is useful to also implement a test method
\begin{Verbatim}
def isEmpty: Boolean
\end{Verbatim}
for sets.
\end{samepage}

\begin{solution}
	\VerbatimInput{\codeDir/intset/IntSet.scala}
\end{solution}

\question
Write an implementation \verb!Integer! of integer numbers. 
The implementation should support all operations of class \verb!Nat!:
\VerbatimInput{\codeDir/Nat.scala}
while adding two methods:
\begin{Verbatim}
def isPositive: Boolean
def negate: Integer
\end{Verbatim}
The first method should return \verb!true! if the number is positive. 
The second method should negate the number. 
Do not use any of Scala’s standard numeric classes in your implementation. 

(\emph{Hint:} There are two possible ways to implement \verb!Integer!. 
One can either make use the existing implementation of \verb!Nat!, 
representing an integer as a natural number and a sign. 
Alternatively one can generalise the given implementation of \verb!Nat! to \verb!Integer!, 
using the three subclasses \verb!Zero! for 0, \verb!Succ! for positive numbers 
and \verb!Pred! for negative numbers.)

\begin{solution}
\VerbatimInput{\codeDir/Integer.scala}
\end{solution}

\section*{Sequences}
\question
Consider the following definitions representing trees of integers. 
These definitions can be seen as an alternative representation of \verb!IntSet!:

\begin{Verbatim}
abstract class IntTree
case object EmptyTree extends IntTree
case class Node(elem: Int, left: IntTree, right: IntTree) extends IntTree
\end{Verbatim}
Complete the following implementations of function \verb!contains! and 
\verb!insert! for \verb!IntTree!’s.
\begin{Verbatim}
def contains(t: IntTree, v: Int): Boolean = t match { ...
        ...
}
def insert(t: IntTree, v: Int): IntTree = t match { ...
        ...
}	
\end{Verbatim}

\begin{solution}
\VerbatimInput{\codeDir/IntTree.scala}	
\end{solution}

\question
As an example of how lists can be processed, consider sorting the elements of a list of 
numbers into ascending order. One simple way to do so is \emph{insertion sort}, 
which works as follows: 
\begin{quotation}
To sort a non-empty list with first element \verb!x! and rest \verb!xs!, 
sort the remainder \verb!xs! and insert the element \verb!x! at the 
right position in the result. Sorting an empty list will yield the empty list. 
\end{quotation}
Expressed as Scala code:
\begin{Verbatim}
def isort(xs: List[Int]): List[Int] =
        if (xs.isEmpty) Nil
        else insert(xs.head, isort(xs.tail))		
\end{Verbatim}
Provide an implementation of the missing function \verb!insert!.

\begin{solution}
\VerbatimInput{\codeDir/Sort.scala}	
\end{solution}

\begin{samepage}
\question
Given the following function, which computes the length of a list:
\begin{Verbatim}
def length: Int = this match {
        case Nil => 0
        case x :: xs => 1 + xs.length
}	
\end{Verbatim}
Design a tail-recursive version of \verb!length!.
\end{samepage}

\begin{solution}
\VerbatimInput{\codeDir/AList.scala}	
\end{solution}

\question
Consider a function which squares all elements of a list and returns a 
list with the results. Complete the following two equivalent definitions of \verb!squareList!:
\begin{Verbatim}
def squareList(xs: List[Int]): List[Int] = xs match {
        case List() => ??
        case y :: ys => ??
}
def squareList(xs: List[Int]): List[Int] =
        xs map ??	
\end{Verbatim}

\begin{solution}
\VerbatimInput{\codeDir/Square.scala}	
\end{solution}

\question
Considering the following methods of class \verb!List!:
\begin{Verbatim}
def filter(p: A => Boolean): List[A] = this match {
    case Nil => this
    case x :: xs => if (p(x)) x :: xs.filter(p) else xs.filter(p)
}

def forall(p: A => Boolean): Boolean =
    isEmpty || (p(head) && (tail forall p))

def exists(p: A => Boolean): Boolean =
    !isEmpty && (p(head) || (tail exists p))	
\end{Verbatim}
Define \verb!forall! and \verb!exists! in terms of \verb!filter!.

\VerbatimInput{\codeDir/List.scala}

\begin{solution}
\ldots	
\end{solution}

\question
Consider the problem of writing a function \verb!flatten!, 
which takes a list of element lists as arguments. The result of \verb!flatten! 
should be the concatenation of all element lists into a single list. 
Here is an implementation of this method in terms of \verb!:\!:
\begin{Verbatim}
def flatten[A](xs: List[List[A]]): List[A] =
        (xs :\ (Nil: List[A])) {(x, xs) => x ::: xs}	
\end{Verbatim}
Consider replacing the body of \verb!flatten! by
\begin{Verbatim}
((Nil: List[A]) /: xs) ((xs, x) => xs ::: x)	
\end{Verbatim}

\begin{solution}
\VerbatimInput{\codeDir/Flat.scala}	
\end{solution}

\question
Complete the missing expressions in the following definitions of some basic 
list-manipulation operations as \emph{fold} operations.
\begin{Verbatim}
def mapFun[A, B](xs: List[A], f: A => B): List[B] =
    (xs :\ List[B]()){ ?? }

def lengthFun[A](xs: List[A]): Int =
    (0 /: xs){ ?? }	
\end{Verbatim}

\begin{solution}
\ldots	
\end{solution}

\question
Given a chess board of size \verb!n!, place \verb!n! 
queens on it so that no queen is in check with another. 
Consider the following code:

\VerbatimInput{\codeDir/Queens.scala}   

Write the function
\begin{Verbatim}
def isSafe(col: Int, queens: List[Int], delta: Int): Boolean	
\end{Verbatim}
which tests whether a queen in the given column \verb!col! is \emph{safe} 
with respect to the queens already placed. Here, \verb!delta! is the difference 
between the row of the queen to be placed and the row of the first queen in the list.

\begin{solution}
\VerbatimInput{\codeDir/IsSafe.scala}
\end{solution}

\question
Define the following function in terms of \verb!for!:
\begin{Verbatim}
def flatten[A](xss: List[List[A]]): List[A] =
  (xss :\ (Nil: List[A])) ((xs, ys) => xs ::: ys)	
\end{Verbatim}

\begin{solution}
\VerbatimInput{\codeDir/Flatten.scala}	
\end{solution}

\section*{Higher-Order functions}

\question
Translate
\begin{Verbatim}
for (b <- books; a <- b.authors if a startsWith "Bird") yield b.title
for (b <- books if (b.title indexOf "Program") >= 0) yield b.title	
\end{Verbatim}
to higher-order functions.

\begin{solution}
\VerbatimInput{\codeDir/Higher.scala}	
\end{solution}

\question % 5.4.6.2 Folding Maybe ES
\begin{parts}
\part
Given the following sum type for modelling optional data:
\begin{Verbatim}
sealed trait Maybe[A]
final case class Full[A](value: A) extends Maybe[A] 
final case class Empty[A]() extends Maybe[A]	
\end{Verbatim}
Implement \verb!fold! for \verb!Maybe!.

\begin{solution}
\begin{Verbatim}
sealed trait Maybe[A] {
  def fold[B](full: A => B, empty: B): B =
    this match {
      case Full(v) => full(v) 
      case Empty() => empty
    } 
}

final case class Full[A](value: A) extends Maybe[A] 
final case class Empty[A]() extends Maybe[A]	
\end{Verbatim}
\end{solution}

\part
Given the following generic sum type:
\begin{Verbatim}
sealed trait Sum[A, B]
final case class Left[A, B](value: A) extends Sum[A, B] 
final case class Right[A, B](value: B) extends Sum[A, B]	
\end{Verbatim}
Implement \verb!fold! for \verb!Sum!.
	
\begin{solution}
\begin{Verbatim}
sealed trait Sum[A, B] {
  def fold[C](left: A => C, right: B => C): C =
    this match {
      case Left(a) => left(a)
      case Right(b) => right(b) 
    }
}
final case class Left[A, B](value: A) extends Sum[A, B] 
final case class Right[A, B](value: B) extends Sum[A, B]	
\end{Verbatim}
\end{solution}

\end{parts}

\question % 5.3.4.1 Tree ES
A binary tree can be defined as follows:
\begin{quotation}
A Tree of type A is a Node with a left and right Tree or a Leaf with an element of type A.
\end{quotation}
\begin{parts}
\part 
Implement this algebraic data type along with a \verb!fold! method.	

\begin{solution}
\begin{Verbatim}
sealed trait Tree[A] {
  def fold[B](node: (B, B) => B, leaf: A => B): B
}

final case class Node[A](left: Tree[A], right: Tree[A]) 
  extends Tree[A] {
    def fold[B](node: (B, B) => B, leaf: A => B): B = 
      node(left.fold(node, leaf), right.fold(node, leaf))
}
    
final case class Leaf[A](value: A) extends Tree[A] {
  def fold[B](node: (B, B) => B, leaf: A => B): B = 
    leaf(value)
}	
\end{Verbatim}
\end{solution}

\part
Using \verb!fold! convert the following \verb!Tree! to a \verb!String!.

\begin{Verbatim}
val tree: Tree[String] = 
  Node(Node(Leaf("To"), Leaf("iterate")),
    Node(Node(Leaf("is"), Leaf("human,")),
      Node(Leaf("to"), Node(Leaf("recurse"), Leaf("divine")))))	
\end{Verbatim}

\begin{solution}
Note it is necessary to instantiate the generic type variable for \verb!fold!. 
Type inference fails in this case.
\begin{Verbatim}
  tree.fold[String]((a, b) => a + " " + b, str => str)	
\end{Verbatim}
\end{solution}

\end{parts}

\end{questions}
\end{document}